\documentclass[12pt,a4paper,oneside]{report}

% --- PACHETE ESENȚIALE ---
\usepackage[utf8]{inputenc}
\usepackage[T1]{fontenc}
\usepackage[romanian]{babel}
\usepackage{graphicx}
\usepackage{geometry}
\geometry{a4paper, top=2.5cm, bottom=2.5cm, left=3cm, right=2.5cm}
\usepackage{times} % Font standard academic
\usepackage{setspace}
\onehalfspacing % Spațiere la 1.5 rânduri
\usepackage{hyperref}
\usepackage{listings}
\usepackage{xcolor}
\usepackage{titlesec}

% --- CONFIGURARE COD SURSĂ (pentru Dart/JS) ---
\definecolor{codegreen}{rgb}{0,0.6,0}
\definecolor{codegray}{rgb}{0.5,0.5,0.5}
\definecolor{codepurple}{rgb}{0.58,0,0.82}
\lstset{
    commentstyle=\color{codegreen},
    keywordstyle=\color{magenta},
    numberstyle=	iny\color{codegray},
    stringstyle=\color{codepurple},
    basicstyle=	tfamily\footnotesize,
    breakatwhitespace=false,         
    breaklines=true,                 
    captionpos=b,                    
    keepspaces=true,                 
    numbers=left,                    
    numbersep=5pt,                  
    showspaces=false,                
    showstringspaces=false,
    showtabs=false,                  
    tabsize=2
}

% --- DATELE LUCRIĂRII ---
	itle{Sistem distribuit pentru managementul productivității și controlul distracțiilor digitale}
\author{Numele Tău}
\date{Iunie 2026}

\begin{document}

% --- PAGINA DE TITLU (Simplificată - va trebui adaptată la formatul UPT) ---
\begin{titlepage}
    \begin{center}
        \large Universitatea Politehnica Timișoara 
        Facultatea de Automatică și Calculatoare 
        Specializarea [Numele Specializării Tale - ex: CTI / IS] 
        \vfill
        \Huge 	extbf{LUCRARE DE LICENȚĂ} 
        \vspace{1cm}
        \Large 	extbf{Sistem distribuit pentru managementul productivității și controlul distracțiilor digitale} 
        \vfill
        \begin{flushleft}
            \large 	extbf{Coordonator științific:} 
            [Titlu și Nume Coordonator]
        \end{flushleft}
        \begin{flushright}
            \large 	extbf{Absolvent:} 
            [Numele Tău]
        \end{flushright}
        \vfill
        \large Timișoara 
        2026
    \end{center}
\end{titlepage}

	ableofcontents

\chapter{Introducere}
\section{Motivația proiectului}
\section{Obiectivele lucrării}

\chapter{Analiza domeniului și tehnologii utilizate}
\section{Stadiul actual al aplicațiilor de productivitate}
\section{Tehnologii de dezvoltare cross-platform: Flutter}
\section{Ecosistemul Google Cloud și Firebase Realtime Database}
\section{Extensii de browser și Google Calendar API}

\chapter{Proiectarea și Arhitectura Sistemului}
\section{Arhitectura de ansamblu a sistemului distribuit}
\section{Modelarea datelor în Firebase}
\section{Mecanisme de sincronizare între dispozitive (Presence System)}

\chapter{Implementarea Soluției}
\section{Aplicația Flutter: Managementul proceselor pe Android și Windows}
\section{Extensia de browser: Filtrarea traficului web și sincronizarea sesiunilor}
\section{Integrarea cu Google Calendar pentru automatizarea focusului}

\chapter{Testare și Validare}
\section{Cazuri de utilizare și testarea funcționalității}
\section{Analiza performanței și consumului de resurse}

\chapter{Concluzii și direcții viitoare}

\bibliographystyle{plain}
\begin{thebibliography}{99}
\bibitem{flutter} Flutter Documentation, https://docs.flutter.dev/
\bibitem{firebase} Firebase Realtime Database, https://firebase.google.com/docs/database
\end{thebibliography}

\end{document}
